%=======================================================================
\section{Selecting the Best Hyper-parameters per Descriptor (Phase 1)}
\label{sec:hyper}

\subsection{The problem with brute force}
Each descriptor has many possible settings. Training a classifier for every combination would take
hours and conflate \textbf{feature quality} with \textbf{model quality}. We want to judge the
features themselves, quickly and fairly.

\subsection{Classifier-free class-separability metrics}
Instead, every candidate configuration is scored in milliseconds with three complementary metrics
that measure how well the resulting feature distributions separate the three classes. All features
are standardised first so the comparison is fair across configurations.

\begin{center}
\begin{tabularx}{\textwidth}{@{}l Y l@{}}
\toprule
\thd{Metric} & \thd{What it measures} & \thd{Direction} \\
\midrule
Fisher Discriminant Ratio (FDR) & Between-class variance $\div$ within-class variance, averaged over
features. Large means class means are far apart relative to their spread. & Higher is better \\
Mutual Information (MI) & Statistical dependency between each feature and the class label (averaged).
Captures non-linear relationships FDR misses; scale-invariant. & Higher is better \\
Davies--Bouldin Index (DBI) & Ratio of within-class scatter to between-class distance in
(PCA-reduced) feature space. Measures cluster compactness \& separation. & Lower is better \\
\bottomrule
\end{tabularx}
\end{center}

\subsection{Combining them into one score}
The three metrics are on different scales, so for each descriptor's set of candidates they are
\textbf{min--max normalised to $[0,1]$}, the DBI is inverted (so that higher is always better), and
the three normalised values are averaged:
\begin{equation*}
\text{score} = \frac{\text{FDR}_{\text{norm}} + \text{MI}_{\text{norm}} + (1 - \text{DBI}_{\text{norm}})}{3}.
\end{equation*}
The configuration with the highest composite score is selected for each descriptor. This composite is
robust: a configuration must do well on \textbf{all three} complementary criteria to win, not just
one.

\subsection{Search space explored}
\begin{center}
\begin{tabularx}{\textwidth}{@{}l Y@{}}
\toprule
\thd{Descriptor} & \thd{Configurations evaluated} \\
\midrule
GLCM & 6 distance sets $\times$ 4 angle sets $\times$ 3 levels $\times$ 2 symmetry $= 144$ \\
LBP  & 4 $P$ $\times$ 3 $R$ $\times$ 3 methods $= 36$ \\
DWT  & 7 wavelets $\times$ 3 levels $= 21$ \\
HOG  & 3 orientations $\times$ 2 cell sizes $\times$ 2 block sizes $\times$ 1 norm $= 12$ \\
\midrule
\textbf{Total} & \textbf{213 configurations, all scored without training a classifier} \\
\bottomrule
\end{tabularx}
\end{center}

\subsection{Selected optimal parameters (results)}
Running Phase~1 on the training set produced the following winners:

\begin{center}
\renewcommand{\arraystretch}{1.2}
\begin{tabularx}{\textwidth}{@{}l Y c c@{}}
\toprule
\thd{Descriptor} & \thd{Selected parameters} & \thd{Score} & \thd{Size} \\
\midrule
\textbf{GLCM} & \glcmParams & \glcmScore & 6 \\
\textbf{LBP}  & \lbpParams  & \lbpScore  & 10 \\
\textbf{DWT}  & \dwtParams  & \dwtScore  & 16 \\
\textbf{HOG}  & \hogParams  & \hogScore  & 1176 \\
\bottomrule
\end{tabularx}
\renewcommand{\arraystretch}{1.0}
\end{center}

Reading these results: \textbf{LBP} achieved a perfect composite score (\lbpScore) on the training
set, indicating its uniform $P{=}8$/$R{=}1$ histograms separate the three classes extremely cleanly;
\textbf{DWT (bior1.3)} and \textbf{HOG} are also very strong. \textbf{GLCM} scores lower
(\glcmScore) --- its six averaged properties carry useful but less individually discriminative
information --- yet it is still retained because it is cheap and adds a complementary texture view to
the fusion. The chosen GLCM offsets $[1,2,4,8]$ deliberately span fine to coarse texture, while a
single angle (0) keeps the descriptor compact.
